% book06.tex --- Book VI of Euclid's Elements: Similar Figures.
%
% All 33 propositions encoded. Book VI applies the Eudoxean proportion
% machinery of Book V to plane figures: triangles, parallelograms, and
% similar polygons. Highlights are VI.2 (the intercept theorem / "basic
% proportionality theorem"), VI.4-VI.8 (similarity criteria), VI.13
% (mean proportional construction), VI.30 (golden section by
% application of areas), and VI.31 (generalised Pythagoras).
%
% Wording follows Heath (1908).

\section{Book VI --- Similar Figures}
\label{sec:book-VI}

\begin{claim}[Proposition VI.1: Triangles and parallelograms with the same height]
\label{prop:VI.1}
Triangles and parallelograms which are under the same height are to
one another as their bases.
\end{claim}
\begin{evidence}[Proof of VI.1]
\label{ev:VI.1}
Repeated application of I.38 generates any multiple of a triangle on
the corresponding multiple of the base; the resulting equimultiples
test (V.5) is exactly the statement of proportionality.  The
parallelogram version follows because each parallelogram is double
its diagonal triangle (I.34, I.41).
\dependson{VI.1}{I.34}
\dependson{VI.1}{I.38}
\dependson{VI.1}{I.41}
\dependson{VI.1}{def:V.5}
\end{evidence}

\begin{claim}[Proposition VI.2: Parallel to a side cuts proportionally]
\label{prop:VI.2}
If a straight line be drawn parallel to one of the sides of a
triangle, it will cut the sides of the triangle proportionally; and
if the sides of the triangle be cut proportionally, the line joining
the points of section will be parallel to the remaining side of the
triangle.
\end{claim}
\begin{evidence}[Proof of VI.2]
\label{ev:VI.2}
Drop a parallel $DE$ from a point $D$ on $AB$ to a point $E$ on $AC$,
parallel to $BC$.  Triangles $\triangle DBC$ and $\triangle DCE$
share the same base $DE$ and lie between the same parallels (with
$BE$, $DC$ as transversals through I.29, I.37), so they are equal in
area.  Applying VI.1 to $\triangle ADE$ versus the equal-area
companion triangles gives $AD : DB = AE : EC$.  The converse runs the
argument in reverse via I.39.
\dependson{VI.2}{I.29}
\dependson{VI.2}{I.37}
\dependson{VI.2}{I.38}
\dependson{VI.2}{I.39}
\dependson{VI.2}{VI.1}
\dependson{VI.2}{V.11}
\end{evidence}

\begin{claim}[Proposition VI.3: Angle bisector cuts opposite side proportionally]
\label{prop:VI.3}
If an angle of a triangle be bisected and the straight line cutting
the angle cut the base also, the segments of the base will have the
same ratio as the remaining sides of the triangle.
\end{claim}
\begin{evidence}[Proof of VI.3]
\label{ev:VI.3}
Through $C$ draw $CE$ parallel to the bisector $AD$ (I.31), meeting
$BA$ produced at $E$.  By alternate angles (I.29) and the bisection
hypothesis, $\angle ACE = \angle AEC$, so $AE = AC$ (I.6).  Applying
VI.2 to $\triangle BCE$ with $AD \parallel CE$: $BD : DC = BA : AE =
BA : AC$.
\dependson{VI.3}{I.6}
\dependson{VI.3}{I.29}
\dependson{VI.3}{I.31}
\dependson{VI.3}{VI.2}
\end{evidence}

\begin{claim}[Proposition VI.4: Equiangular triangles have proportional sides]
\label{prop:VI.4}
In equiangular triangles the sides about the equal angles are
proportional, and those are corresponding sides which subtend the
equal angles.
\end{claim}
\begin{evidence}[Proof of VI.4]
\label{ev:VI.4}
Lay the equiangular triangles side by side so that one pair of equal
angles coincides; the remaining vertices and bases yield a
parallelogram by I.28.  Apply VI.2 to the new figure to derive the
proportionality of the sides.
\dependson{VI.4}{I.28}
\dependson{VI.4}{I.32}
\dependson{VI.4}{VI.2}
\end{evidence}

\begin{claim}[Proposition VI.5: SSS-similarity criterion]
\label{prop:VI.5}
If two triangles have their sides proportional, the triangles will be
equiangular and will have those angles equal which the corresponding
sides subtend.
\end{claim}
\begin{evidence}[Proof of VI.5]
\label{ev:VI.5}
Construct on the second triangle's base a triangle equiangular with
the first (I.23); by VI.4 its other sides are determined by the
proportion, and by I.8 (SSS) it coincides with the second triangle.
Therefore the second triangle is equiangular with the first.
\dependson{VI.5}{I.8}
\dependson{VI.5}{I.23}
\dependson{VI.5}{VI.4}
\end{evidence}

\begin{claim}[Proposition VI.6: SAS-similarity criterion]
\label{prop:VI.6}
If two triangles have one angle equal to one angle and the sides
about the equal angles proportional, the triangles will be
equiangular and will have those angles equal which the corresponding
sides subtend.
\end{claim}
\begin{evidence}[Proof of VI.6]
\label{ev:VI.6}
Same scheme as VI.5: extend one triangle so as to match the second
on the equal-angle pair (I.23), apply VI.4 to deduce the missing
side, then I.4 (SAS) for congruence of the auxiliary triangle with
the second.
\dependson{VI.6}{I.4}
\dependson{VI.6}{I.23}
\dependson{VI.6}{VI.4}
\dependson{VI.6}{VI.5}
\end{evidence}

\begin{claim}[Proposition VI.7: Mixed SAS-similarity criterion]
\label{prop:VI.7}
If two triangles have one angle equal to one angle, the sides about
other angles proportional, and the remaining angles either both less
or both not less than a right angle, the triangles will be
equiangular and will have those angles equal about which the sides
are proportional.
\end{claim}
\begin{evidence}[Proof of VI.7]
\label{ev:VI.7}
Construct, at the vertex of the angle whose sides are proportional,
an angle equal to the corresponding angle in the second triangle
(I.23).  The resulting auxiliary triangle agrees with the first in
two angles (and hence all three, by I.32) and shares a side with the
second; the constraint on the remaining angle being acute or obtuse
ensures the construction is non-ambiguous (essentially eliminates
the SSA failure case).
\dependson{VI.7}{I.23}
\dependson{VI.7}{I.32}
\dependson{VI.7}{VI.4}
\dependson{VI.7}{VI.6}
\end{evidence}

\begin{claim}[Proposition VI.8: Right triangle altitude similarity]
\label{prop:VI.8}
If in a right-angled triangle a perpendicular be drawn from the right
angle to the base, the triangles adjoining the perpendicular are
similar both to the whole and to one another.
\end{claim}
\begin{evidence}[Proof of VI.8]
\label{ev:VI.8}
The two sub-triangles each share an angle with the original (the
non-right angle at $B$ or $C$) and both have a right angle (at the
foot of the altitude and at the apex), so they are equiangular with
the original by I.32, hence similar by VI.4.  By transitivity (V.11)
they are similar to each other.
\dependson{VI.8}{I.32}
\dependson{VI.8}{VI.4}
\dependson{VI.8}{V.11}
\end{evidence}

\begin{claim}[Proposition VI.9: Cut off any part of a given segment]
\label{prop:VI.9}
From a given straight line to cut off a prescribed part.
\end{claim}
\begin{evidence}[Proof of VI.9]
\label{ev:VI.9}
Lay off a separate transversal containing the prescribed number of
equal units (I.3 repeatedly).  Join the far ends and draw parallels
to that join through each unit division (I.31).  By VI.2 these
parallels mark off equal fractions on the given line.
\dependson{VI.9}{I.3}
\dependson{VI.9}{I.31}
\dependson{VI.9}{VI.2}
\end{evidence}

\begin{claim}[Proposition VI.10: Divide a segment in a given ratio]
\label{prop:VI.10}
To cut a given uncut straight line similarly to a given cut straight
line.
\end{claim}
\begin{evidence}[Proof of VI.10]
\label{ev:VI.10}
Lay the cut transversal alongside the given line meeting at one
endpoint; join the far ends and draw parallels to that join through
each cut point (I.31).  By VI.2 the parallels reproduce the same
ratios on the given line.
\dependson{VI.10}{I.31}
\dependson{VI.10}{VI.2}
\end{evidence}

\begin{claim}[Proposition VI.11: Third proportional]
\label{prop:VI.11}
To two given straight lines to find a third proportional.
\end{claim}
\begin{evidence}[Proof of VI.11]
\label{ev:VI.11}
Place the two given segments on lines making an angle.  Mark off the
second segment beyond the first; draw a parallel to the closing
segment through the far endpoint (I.31).  By VI.2 the new mark-off
is the required third proportional.
\dependson{VI.11}{I.31}
\dependson{VI.11}{VI.2}
\end{evidence}

\begin{claim}[Proposition VI.12: Fourth proportional]
\label{prop:VI.12}
To three given straight lines to find a fourth proportional.
\end{claim}
\begin{evidence}[Proof of VI.12]
\label{ev:VI.12}
Same construction as VI.11 but with three input segments: lay two on
one transversal and one on the other, then drop a parallel from the
last point.  VI.2 yields the fourth proportional.
\dependson{VI.12}{I.31}
\dependson{VI.12}{VI.2}
\end{evidence}

\begin{claim}[Proposition VI.13: Mean proportional]
\label{prop:VI.13}
To two given straight lines to find a mean proportional.
\end{claim}
\begin{evidence}[Proof of VI.13]
\label{ev:VI.13}
Lay the two segments end-to-end as $AB$, $BC$ on a single line; on
$AC$ as diameter describe a semicircle (Postulate 3, III.31 implicit).
Erect the perpendicular $BD$ at $B$ to the diameter; then $\triangle
ABD$, $\triangle BDC$, and $\triangle ABC$ are similar by VI.8, so
$AB : BD = BD : BC$, making $BD$ the required mean proportional.
\dependson{VI.13}{I.11}
\dependson{VI.13}{III.31}
\dependson{VI.13}{VI.8}
\dependson{VI.13}{post:3}
\end{evidence}

\begin{claim}[Proposition VI.14: Equal parallelograms have reciprocally proportional sides]
\label{prop:VI.14}
In equal and equiangular parallelograms the sides about the equal
angles are reciprocally proportional; and equiangular parallelograms
in which the sides about the equal angles are reciprocally
proportional are equal.
\end{claim}
\begin{evidence}[Proof of VI.14]
\label{ev:VI.14}
Lay the two parallelograms so the equal angles coincide; their union
forms a third parallelogram whose diagonal includes the original
common-angle vertex.  By VI.1 the ratios of areas equal the ratios of
adjacent sides; equality of the original areas forces the
reciprocal-proportion relation.  Converse runs the same way.
\dependson{VI.14}{VI.1}
\dependson{VI.14}{V.11}
\end{evidence}

\begin{claim}[Proposition VI.15: Equal triangles with one common angle have reciprocally proportional sides]
\label{prop:VI.15}
In equal triangles which have one angle equal to one angle the sides
about the equal angles are reciprocally proportional; and those
triangles which have one angle equal to one angle, and in which the
sides about the equal angles are reciprocally proportional, are equal.
\end{claim}
\begin{evidence}[Proof of VI.15]
\label{ev:VI.15}
Same scheme as VI.14 applied to triangles (each is half a
parallelogram, so the areas-of-parallelograms result transfers).
\dependson{VI.15}{I.41}
\dependson{VI.15}{VI.14}
\end{evidence}

\begin{claim}[Proposition VI.16: Equal rectangles iff proportional sides]
\label{prop:VI.16}
If four straight lines be proportional, the rectangle contained by the
extremes is equal to the rectangle contained by the means; and if the
rectangle contained by the extremes be equal to the rectangle
contained by the means, the four straight lines will be proportional.
\end{claim}
\begin{evidence}[Proof of VI.16]
\label{ev:VI.16}
If $a : b = c : d$, the rectangles $ad$ and $bc$ are equiangular
(all right-angled) and have sides reciprocally proportional, so by
VI.14 they are equal.  Conversely from $ad = bc$ and VI.14 (its
converse) we recover the proportion.
\dependson{VI.16}{VI.14}
\end{evidence}

\begin{claim}[Proposition VI.17: Mean proportional iff equal squares]
\label{prop:VI.17}
If three straight lines be proportional, the rectangle contained by
the extremes is equal to the square on the mean; and if the rectangle
contained by the extremes be equal to the square on the mean, the
three straight lines will be proportional.
\end{claim}
\begin{evidence}[Proof of VI.17]
\label{ev:VI.17}
Specialise VI.16 to $b = c$.
\dependson{VI.17}{VI.16}
\end{evidence}

\begin{claim}[Proposition VI.18: Construct a polygon similar to a given polygon]
\label{prop:VI.18}
On a given straight line to describe a rectilineal figure similar
and similarly situated to a given rectilineal figure.
\end{claim}
\begin{evidence}[Proof of VI.18]
\label{ev:VI.18}
Triangulate the given polygon by joining a vertex to all others.  On
the new base copy each angle (I.23) and use VI.4 to fix the side
ratios; assemble the triangles into the similar polygon.
\dependson{VI.18}{I.23}
\dependson{VI.18}{VI.4}
\end{evidence}

\begin{claim}[Proposition VI.19: Similar triangles are as the squares on corresponding sides]
\label{prop:VI.19}
Similar triangles are to one another in the duplicate ratio of the
corresponding sides.
\end{claim}
\begin{evidence}[Proof of VI.19]
\label{ev:VI.19}
Let $\triangle ABC \sim \triangle DEF$ with side ratio $k = BC / EF$.
Construct $G$ on $BC$ so that $BG = EF \cdot k$ (i.e.\ a third
proportional, VI.11).  By VI.1 the area ratio is $BG : EF = k$ along
one dimension and the side ratio $k$ along the perpendicular, giving
total area ratio $k^2$ in the sense of Definition V.9 (duplicate
ratio).
\dependson{VI.19}{VI.1}
\dependson{VI.19}{VI.11}
\dependson{VI.19}{def:V.9}
\end{evidence}

\begin{claim}[Proposition VI.20: Similar polygons are as the squares on corresponding sides]
\label{prop:VI.20}
Similar polygons are divided into similar triangles, equal in
multitude and in the same ratio as the wholes; and the polygon has
to the polygon a ratio duplicate of that which the corresponding side
has to the corresponding side.
\end{claim}
\begin{evidence}[Proof of VI.20]
\label{ev:VI.20}
Decompose into similar triangles by joining one vertex to all others;
apply VI.19 to each triangle and sum via V.12.
\dependson{VI.20}{VI.18}
\dependson{VI.20}{VI.19}
\dependson{VI.20}{V.12}
\end{evidence}

\begin{claim}[Proposition VI.21: Figures similar to the same are similar]
\label{prop:VI.21}
Figures which are similar to the same rectilineal figure are also
similar to one another.
\end{claim}
\begin{evidence}[Proof of VI.21]
\label{ev:VI.21}
Equiangularity transfers transitively, and the side ratios compose
transitively by V.11.
\dependson{VI.21}{V.11}
\dependson{VI.21}{VI.4}
\end{evidence}

\begin{claim}[Proposition VI.22: Proportionality of figures on proportional sides]
\label{prop:VI.22}
If four straight lines be proportional, the rectilineal figures
similar and similarly described upon them will also be proportional;
and if the rectilineal figures similar and similarly described upon
them be proportional, the straight lines will themselves also be
proportional.
\end{claim}
\begin{evidence}[Proof of VI.22]
\label{ev:VI.22}
Apply VI.20: figures similar on proportional sides have areas in the
duplicate ratio of the sides.  Equality of those duplicate ratios is
equivalent to equality of the original side ratios.
\dependson{VI.22}{VI.20}
\dependson{VI.22}{def:V.9}
\end{evidence}

\begin{claim}[Proposition VI.23: Ratio of parallelograms is compound of side ratios]
\label{prop:VI.23}
Equiangular parallelograms have to one another the ratio compounded
of the ratios of their sides.
\end{claim}
\begin{evidence}[Proof of VI.23]
\label{ev:VI.23}
Place the two parallelograms so the equal angles share a vertex; the
resulting figure can be split by lines parallel to the sides into a
rectangle whose dimensions are the four sides.  Two applications of
VI.1 give the compounded ratio.
\dependson{VI.23}{VI.1}
\dependson{VI.23}{def:V.10}
\end{evidence}

\begin{claim}[Proposition VI.24: Parallelograms about the diagonal are similar to the whole]
\label{prop:VI.24}
In any parallelogram the parallelograms about the diameter are
similar both to the whole and to one another.
\end{claim}
\begin{evidence}[Proof of VI.24]
\label{ev:VI.24}
Lines drawn parallel to the sides through a point on the diagonal
make the inner parallelograms equiangular with the whole (I.29) and
with corresponding sides cut in the same ratio (VI.2); hence similar
(VI.4).
\dependson{VI.24}{I.29}
\dependson{VI.24}{VI.2}
\dependson{VI.24}{VI.4}
\end{evidence}

\begin{claim}[Proposition VI.25: Construct a figure similar to one and equal to another]
\label{prop:VI.25}
To construct one and the same figure similar to a given rectilineal
figure and equal to another given rectilineal figure.
\end{claim}
\begin{evidence}[Proof of VI.25]
\label{ev:VI.25}
Reduce both given figures to rectangles on a common base (I.44);
the side opposite the common base measures each figure's area.  Take
the mean proportional (VI.13) of those two opposite sides; build on
that mean a figure similar to the first via VI.18.  By VI.20 the
constructed figure has the required area.
\dependson{VI.25}{I.44}
\dependson{VI.25}{VI.13}
\dependson{VI.25}{VI.18}
\dependson{VI.25}{VI.20}
\end{evidence}

\begin{claim}[Proposition VI.26: Inner parallelogram about the diagonal must share a corner]
\label{prop:VI.26}
If from a parallelogram there be taken away a parallelogram similar
and similarly situated to the whole and having a common angle with
it, it is about the same diameter with the whole.
\end{claim}
\begin{evidence}[Proof of VI.26]
\label{ev:VI.26}
Reductio: if the inner parallelogram is not about the diagonal, then
VI.24 forces a similar parallelogram on the actual diagonal that
differs from the given one; matching corresponding sides via the
similarity contradicts the assumed common-angle configuration.
\dependson{VI.26}{VI.24}
\dependson{VI.26}{V.9}
\end{evidence}

\begin{claim}[Proposition VI.27: Maximum-area parallelogram with deficiency]
\label{prop:VI.27}
Of all parallelograms applied to the same straight line and deficient
by parallelogrammic figures similar and similarly situated to that
described upon the half of the straight line, the greatest is that
which is applied to the half and is similar to the deficient figure.
\end{claim}
\begin{evidence}[Proof of VI.27]
\label{ev:VI.27}
Comparison of areas via VI.20 and VI.24: the application to the half
exhausts the bound, while any other application produces a smaller
parallelogram by a square-on-the-deviation.
\dependson{VI.27}{VI.20}
\dependson{VI.27}{VI.24}
\dependson{VI.27}{VI.26}
\end{evidence}

\begin{claim}[Proposition VI.28: Apply a parallelogram with deficiency]
\label{prop:VI.28}
To a given straight line to apply a parallelogram equal to a given
rectilineal figure and deficient by a parallelogrammic figure similar
to a given one; thus the given rectilineal figure must not be greater
than the parallelogram described on the half of the straight line and
similar to the defect.
\end{claim}
\begin{evidence}[Proof of VI.28]
\label{ev:VI.28}
Construct the application by VI.25 to match the given figure, then
verify the bound via VI.27.  The construction effectively solves a
quadratic equation in geometric form.
\dependson{VI.28}{VI.25}
\dependson{VI.28}{VI.27}
\end{evidence}

\begin{claim}[Proposition VI.29: Apply a parallelogram with excess]
\label{prop:VI.29}
To a given straight line to apply a parallelogram equal to a given
rectilineal figure and exceeding by a parallelogrammic figure similar
to a given one.
\end{claim}
\begin{evidence}[Proof of VI.29]
\label{ev:VI.29}
Dual to VI.28; solves the geometric quadratic with the opposite sign.
\dependson{VI.29}{VI.25}
\dependson{VI.29}{VI.28}
\end{evidence}

\begin{claim}[Proposition VI.30: Golden section by application of areas]
\label{prop:VI.30}
To cut a given finite straight line in extreme and mean ratio.
\end{claim}
\begin{evidence}[Proof of VI.30]
\label{ev:VI.30}
Apply to the given line a parallelogram equal to the square on it,
exceeding by a square (VI.29).  The exceeding square's side $x$
satisfies $x(a+x) = a^2$, equivalently $a : x = x : (a - x)$ in the
language of II.11 -- the golden section.
\dependson{VI.30}{II.11}
\dependson{VI.30}{VI.17}
\dependson{VI.30}{VI.29}
\end{evidence}

\begin{claim}[Proposition VI.31: Generalised Pythagoras]
\label{prop:VI.31}
In right-angled triangles the figure on the side subtending the right
angle is equal to the similar and similarly described figures on the
sides containing the right angle.
\end{claim}
\begin{evidence}[Proof of VI.31]
\label{ev:VI.31}
By VI.8 the altitude from the right angle cuts the hypotenuse into
segments such that each leg is the mean proportional between the
hypotenuse and the adjacent segment.  Applying VI.20 (areas of
similar figures are in the duplicate ratio of corresponding sides)
gives each leg-figure equal to its adjacent piece of the
hypotenuse-figure.  Summing the two pieces by V.24 yields the
hypotenuse-figure.
\dependson{VI.31}{VI.8}
\dependson{VI.31}{VI.19}
\dependson{VI.31}{VI.20}
\dependson{VI.31}{V.24}
\end{evidence}

\begin{claim}[Proposition VI.32: Similarly placed triangles share a vertex angle]
\label{prop:VI.32}
If two triangles having two sides proportional to two sides be placed
together at one angle so that their corresponding sides are also
parallel, the remaining sides of the triangles will be in a straight
line.
\end{claim}
\begin{evidence}[Proof of VI.32]
\label{ev:VI.32}
Equal angles and proportional sides (VI.6) force the triangles to
share the same straight-line continuation at the meeting vertex (I.14).
\dependson{VI.32}{I.14}
\dependson{VI.32}{I.29}
\dependson{VI.32}{VI.6}
\end{evidence}

\begin{claim}[Proposition VI.33: Inscribed angles are as their arcs]
\label{prop:VI.33}
In equal circles angles have the same ratio as the circumferences on
which they stand, whether they stand at the centres or at the
circumferences.
\end{claim}
\begin{evidence}[Proof of VI.33]
\label{ev:VI.33}
Use III.27 (equal arcs subtend equal central angles in equal circles)
to set up an equimultiples test: for any positive integers $m$, $n$,
$m$ copies of one angle correspond to $m$ copies of its arc, and the
order of $m \alpha$ versus $n \beta$ matches the order of $m \cdot
\mathrm{arc}(\alpha)$ versus $n \cdot \mathrm{arc}(\beta)$.  This is
exactly Definition V.5 for $\alpha : \beta = \mathrm{arc}(\alpha) :
\mathrm{arc}(\beta)$.  The inscribed-angle case follows from III.20
(inscribed angle is half the central).
\dependson{VI.33}{III.20}
\dependson{VI.33}{III.27}
\dependson{VI.33}{def:V.5}
\end{evidence}
